% Options for packages loaded elsewhere
\PassOptionsToPackage{unicode}{hyperref}
\PassOptionsToPackage{hyphens}{url}
\PassOptionsToPackage{dvipsnames,svgnames,x11names}{xcolor}
%
\documentclass[
  11pt,
  a4paper,
]{article}
\usepackage{amsmath,amssymb}
\usepackage{iftex}
\ifPDFTeX
  \usepackage[T1]{fontenc}
  \usepackage[utf8]{inputenc}
  \usepackage{textcomp} % provide euro and other symbols
\else % if luatex or xetex
  \usepackage{unicode-math} % this also loads fontspec
  \defaultfontfeatures{Scale=MatchLowercase}
  \defaultfontfeatures[\rmfamily]{Ligatures=TeX,Scale=1}
\fi
\usepackage{lmodern}
\ifPDFTeX\else
  % xetex/luatex font selection
\fi
% Use upquote if available, for straight quotes in verbatim environments
\IfFileExists{upquote.sty}{\usepackage{upquote}}{}
\IfFileExists{microtype.sty}{% use microtype if available
  \usepackage[]{microtype}
  \UseMicrotypeSet[protrusion]{basicmath} % disable protrusion for tt fonts
}{}
\makeatletter
\@ifundefined{KOMAClassName}{% if non-KOMA class
  \IfFileExists{parskip.sty}{%
    \usepackage{parskip}
  }{% else
    \setlength{\parindent}{0pt}
    \setlength{\parskip}{6pt plus 2pt minus 1pt}}
}{% if KOMA class
  \KOMAoptions{parskip=half}}
\makeatother
\usepackage{xcolor}
\usepackage[margin=23mm]{geometry}
\usepackage{color}
\usepackage{fancyvrb}
\newcommand{\VerbBar}{|}
\newcommand{\VERB}{\Verb[commandchars=\\\{\}]}
\DefineVerbatimEnvironment{Highlighting}{Verbatim}{commandchars=\\\{\}}
% Add ',fontsize=\small' for more characters per line
\newenvironment{Shaded}{}{}
\newcommand{\AlertTok}[1]{\textcolor[rgb]{1.00,0.00,0.00}{\textbf{#1}}}
\newcommand{\AnnotationTok}[1]{\textcolor[rgb]{0.38,0.63,0.69}{\textbf{\textit{#1}}}}
\newcommand{\AttributeTok}[1]{\textcolor[rgb]{0.49,0.56,0.16}{#1}}
\newcommand{\BaseNTok}[1]{\textcolor[rgb]{0.25,0.63,0.44}{#1}}
\newcommand{\BuiltInTok}[1]{\textcolor[rgb]{0.00,0.50,0.00}{#1}}
\newcommand{\CharTok}[1]{\textcolor[rgb]{0.25,0.44,0.63}{#1}}
\newcommand{\CommentTok}[1]{\textcolor[rgb]{0.38,0.63,0.69}{\textit{#1}}}
\newcommand{\CommentVarTok}[1]{\textcolor[rgb]{0.38,0.63,0.69}{\textbf{\textit{#1}}}}
\newcommand{\ConstantTok}[1]{\textcolor[rgb]{0.53,0.00,0.00}{#1}}
\newcommand{\ControlFlowTok}[1]{\textcolor[rgb]{0.00,0.44,0.13}{\textbf{#1}}}
\newcommand{\DataTypeTok}[1]{\textcolor[rgb]{0.56,0.13,0.00}{#1}}
\newcommand{\DecValTok}[1]{\textcolor[rgb]{0.25,0.63,0.44}{#1}}
\newcommand{\DocumentationTok}[1]{\textcolor[rgb]{0.73,0.13,0.13}{\textit{#1}}}
\newcommand{\ErrorTok}[1]{\textcolor[rgb]{1.00,0.00,0.00}{\textbf{#1}}}
\newcommand{\ExtensionTok}[1]{#1}
\newcommand{\FloatTok}[1]{\textcolor[rgb]{0.25,0.63,0.44}{#1}}
\newcommand{\FunctionTok}[1]{\textcolor[rgb]{0.02,0.16,0.49}{#1}}
\newcommand{\ImportTok}[1]{\textcolor[rgb]{0.00,0.50,0.00}{\textbf{#1}}}
\newcommand{\InformationTok}[1]{\textcolor[rgb]{0.38,0.63,0.69}{\textbf{\textit{#1}}}}
\newcommand{\KeywordTok}[1]{\textcolor[rgb]{0.00,0.44,0.13}{\textbf{#1}}}
\newcommand{\NormalTok}[1]{#1}
\newcommand{\OperatorTok}[1]{\textcolor[rgb]{0.40,0.40,0.40}{#1}}
\newcommand{\OtherTok}[1]{\textcolor[rgb]{0.00,0.44,0.13}{#1}}
\newcommand{\PreprocessorTok}[1]{\textcolor[rgb]{0.74,0.48,0.00}{#1}}
\newcommand{\RegionMarkerTok}[1]{#1}
\newcommand{\SpecialCharTok}[1]{\textcolor[rgb]{0.25,0.44,0.63}{#1}}
\newcommand{\SpecialStringTok}[1]{\textcolor[rgb]{0.73,0.40,0.53}{#1}}
\newcommand{\StringTok}[1]{\textcolor[rgb]{0.25,0.44,0.63}{#1}}
\newcommand{\VariableTok}[1]{\textcolor[rgb]{0.10,0.09,0.49}{#1}}
\newcommand{\VerbatimStringTok}[1]{\textcolor[rgb]{0.25,0.44,0.63}{#1}}
\newcommand{\WarningTok}[1]{\textcolor[rgb]{0.38,0.63,0.69}{\textbf{\textit{#1}}}}
\usepackage{longtable,booktabs,array}
\usepackage{calc} % for calculating minipage widths
% Correct order of tables after \paragraph or \subparagraph
\usepackage{etoolbox}
\makeatletter
\patchcmd\longtable{\par}{\if@noskipsec\mbox{}\fi\par}{}{}
\makeatother
% Allow footnotes in longtable head/foot
\IfFileExists{footnotehyper.sty}{\usepackage{footnotehyper}}{\usepackage{footnote}}
\makesavenoteenv{longtable}
\setlength{\emergencystretch}{3em} % prevent overfull lines
\providecommand{\tightlist}{%
  \setlength{\itemsep}{0pt}\setlength{\parskip}{0pt}}
\setcounter{secnumdepth}{-\maxdimen} % remove section numbering
\ifLuaTeX
\usepackage[bidi=basic]{babel}
\else
\usepackage[bidi=default]{babel}
\fi
\babelprovide[main,import]{french}
% get rid of language-specific shorthands (see #6817):
\let\LanguageShortHands\languageshorthands
\def\languageshorthands#1{}
\usepackage[dvipsnames]{xcolor}
\usepackage{amsmath,amssymb}
\usepackage{microtype}
\setlength{\emergencystretch}{3em}
\setlength{\parskip}{5pt}
\ifLuaTeX
  \usepackage{selnolig}  % disable illegal ligatures
\fi
\IfFileExists{bookmark.sty}{\usepackage{bookmark}}{\usepackage{hyperref}}
\IfFileExists{xurl.sty}{\usepackage{xurl}}{} % add URL line breaks if available
\urlstyle{same}
\hypersetup{
  pdftitle={Amplification arithmétique - V6.55},
  pdfauthor={Dossier de Vincent P. P. NOËL - rédaction assistée},
  pdflang={fr},
  colorlinks=true,
  linkcolor={NavyBlue},
  filecolor={Maroon},
  citecolor={Blue},
  urlcolor={NavyBlue},
  pdfcreator={LaTeX via pandoc}}

\title{Amplification arithmétique - V6.55}
\usepackage{etoolbox}
\makeatletter
\providecommand{\subtitle}[1]{% add subtitle to \maketitle
  \apptocmd{\@title}{\par {\large #1 \par}}{}{}
}
\makeatother
\subtitle{Restitution exacte, petits conducteurs contrôlés et dette
restante}
\author{Dossier de Vincent P. P. NOËL - rédaction assistée}
\date{14 octobre 2026}

\begin{document}
\maketitle

\hypertarget{ce-que-lessai-a-effectivement-produit}{%
\section{Ce que l'essai a effectivement
produit}\label{ce-que-lessai-a-effectivement-produit}}

\textbf{Type de document :} note de recherche avec preuves auxiliaires,
confrontation constructive et critique, programmes et diagnostics
reproductibles. Version V6.55, qui prolonge le dossier V6.54 ; cette
tentative n'établit aucune nouvelle borne sur l'énergie totale et ne
démontre pas RH. Aucune priorité historique n'est revendiquée.

L'amplification est essayée ici sur les entiers avec le poids exact du
dossier. On construit une famille de traductions arithmétiques, puis on
demande si son énergie moyenne contrôle la somme initiale. Trois
éléments sont démontrés :

\begin{enumerate}
\def\labelenumi{\arabic{enumi}.}
\tightlist
\item
  Les paramètres carrés permettent une \textbf{restitution exacte}, par
  différence entre deux échelles, sans erreur de copie.
\item
  La dette des corrélations se regroupe selon un indice unique \(q\), le
  conducteur d'un caractère quadratique.
\item
  Toute la partie \(q\le H\) est majorée à l'échelle \(HX^2\), à des
  logarithmes près, pour toutes les échelles. Le reste porte sur
  \(q>H\).
\end{enumerate}

\begin{longtable}[]{@{}
  >{\raggedright\arraybackslash}p{(\columnwidth - 4\tabcolsep) * \real{0.3333}}
  >{\raggedright\arraybackslash}p{(\columnwidth - 4\tabcolsep) * \real{0.3333}}
  >{\raggedright\arraybackslash}p{(\columnwidth - 4\tabcolsep) * \real{0.3333}}@{}}
\toprule\noalign{}
\begin{minipage}[b]{\linewidth}\raggedright
Opération
\end{minipage} & \begin{minipage}[b]{\linewidth}\raggedright
Résultat
\end{minipage} & \begin{minipage}[b]{\linewidth}\raggedright
Portée
\end{minipage} \\
\midrule\noalign{}
\endhead
\bottomrule\noalign{}
\endlastfoot
Retrouver la cible dans la famille & Identité exacte et restitution
stable & Démontré \\
Contrôler la diagonale et les petits conducteurs & Bornes explicites,
sans RH & Démontré \\
Contrôler la dette à grand conducteur & Borne supérieure signée uniforme
& Ouvert \\
En déduire un demi-plan sans zéros & \(\Re s>3/4\) si la borne restante
est prouvée & Implication conditionnelle \\
Atteindre la ligne \(1/2\) & Contrôle plus fort de l'énergie initiale &
Ouvert \\
\end{longtable}

\textbf{Verdict.} Le mécanisme de copie fonctionne. Le grand crible
ordinaire ne fournit pas l'annulation nécessaire. Une partie croissante
de la dette est maintenant contrôlée analytiquement ; cela ne constitue
pas un gain d'exposant sur la somme complète. Le verrou restant est
localisé, pas franchi.

Le résultat conditionnel recherché dans cet essai est volontairement
plus faible que RH. Prouver sa dernière estimation constituerait un
changement quantitatif global dans cette construction. Ajouter des
copies ou augmenter un plafond de calcul ne le remplace pas.

\newpage

\hypertarget{la-source-exacte-et-ses-bornes-uxe9luxe9mentaires}{%
\section{1. La source exacte et ses bornes
élémentaires}\label{la-source-exacte-et-ses-bornes-uxe9luxe9mentaires}}

Écrivons \(\Lambda(p^k)=\log p\), et \(\Lambda(n)=0\) sinon. Posons \[
\psi(x)=\sum_{n\le x}\Lambda(n),\qquad
E(x)=\psi(x)-x,\qquad c=1+\gamma_E.
\] Pour \(x\ge1\), avec \(N=\lfloor x\rfloor\), définissons \[
A(x)=E(x)+c,\qquad
B(x)=\log(N!)-x h_N+cN,\qquad
h_N=\sum_{j=1}^N\frac1j.
\] Nous prolongeons \(A\) et \(B\) par zéro sur \(x<1\). Cette
convention s'appliquera à toutes les fonctions et dilatations
ci-dessous.

\textbf{Identité de diviseurs et inversion.} Pour tout \(x\ge1\), \[
\sum_{k\le x}A(x/k)=B(x),\qquad
A(x)=\sum_{k\le x}\mu(k)B(x/k).
\] En effet, \(\sum_{d\mid n}\Lambda(d)=\log n\) entraîne
\(\sum_{k\le x}\psi(x/k)=\log(N!)\). Les deux autres termes donnent
\(-xh_N+cN\). Pour inverser, regrouper les sommes finies par \(j=k\ell\)
et utiliser \(\sum_{k\mid j}\mu(k)=\mathbf1_{j=1}\). Aucune limite
n'intervient.

\textbf{Borne explicite de la source.} Si \(x=N+\theta\),
\(0\le\theta<1\), les restes de Stirling et du développement harmonique
donnent \[
\begin{aligned}
B(x)={}&(\tfrac12-\theta)\log N
 +\tfrac12\log(2\pi)-\tfrac12-\gamma_E\theta\\
 &-\frac{\theta}{2N}+r_N+(N+\theta)e_N,\\
&0<r_N<\frac1{12N},\qquad 0<e_N<\frac1{12N^2}.
\end{aligned}
\] Le reste après \((1/2-\theta)\log N\) est compris entre \(-3/2\) et
\(3/4\). Il en résulte la borne commode \[
|B(x)|\le\tfrac12\log x+2\qquad(x\ge1).
\]

\textbf{Deux coûts utiles.} La comparaison avec les fonctions
décroissantes de \(t\) \(\tfrac12\log(x/t)+2\) et son carré, sur
\(0<t\le x\), donne \[
\sum_{n\le x}|B(x/n)|\le\frac52 x,
\qquad
\sum_{n\le x}|B(x/n)|^2\le\frac{13}{2}x.
\] Par exemple, la seconde somme est au plus
\(\int_0^x(\frac12\log(x/t)+2)^2dt=13x/2\). La diagonale de l'énergie
\(J_A(X)=\int_1^XA(x)^2dx\) vérifie donc \[
D(X)\le\frac{13}{4}(X^2-1).
\] Cela améliore la constante auxiliaire \(13/2\) précédemment annoncée
; cela n'améliore pas la borne de l'énergie totale. Les bornitudes
quadratiques de \(J_A\) et de \(J=\int_1^XE(x)^2dx\) sont équivalentes
par \((a+b)^2\le2a^2+2b^2\).

\newpage

\hypertarget{construire-les-copies-et-les-restituer}{%
\section{2. Construire les copies et les
restituer}\label{construire-les-copies-et-les-restituer}}

Le symbole de Jacobi \((u/n)\) prend les valeurs \(-1,0,1\) pour un
dénominateur impair positif ; \((u/1)=1\). Pour \(n\) sans carré, c'est
un caractère quadratique, nul si \(u\) et \(n\) ont un facteur premier
commun.

Définissons la composante impaire et la famille \[
A_o(x)=\sum_{\substack{n\le x\\n\text{ impair}}}\mu(n)B(x/n),\qquad
F_u(x)=\sum_{\substack{n\le x\\n\text{ impair}}}
\mu(n)\left(\frac un\right)B(x/n).
\] Les termes non sans carré sont automatiquement nuls. En séparant les
indices pairs dans l'inversion de Möbius, \[
\boxed{A(x)=A_o(x)-A_o(x/2).}
\]

\textbf{Lemme de restitution.} Pour tout premier impair \(p\), \[
\boxed{A_o(x)=F_{p^2}(x)-F_{p^2}(x/p).}
\] En effet, \((p^2/n)=\mathbf1_{(n,p)=1}\) pour les indices retenus. La
somme \(F_{p^2}\) est donc \(A_o\) privée des multiples de \(p\). Ces
multiples sont \(n=pm\), avec \((m,p)=1\) ; leur contribution est
\(-F_{p^2}(x/p)\), par \(\mu(pm)=-\mu(m)\). C'est une identité finie,
pas une approximation.

\textbf{Opérateur réutilisable.} Posons \(L_pf(x)=f(x/p)\). Pour \[
\mathcal N_T(f)=\sup_{1\le Y\le T}
\frac1Y\left(\int_1^Y|f(x)|^2dx\right)^{1/2},
\] le changement de variable donne
\(\mathcal N_T(L_pf)\le p^{-1/2}\mathcal N_T(f)\). La restitution
s'écrit \(A_o=(I-L_p)F_{p^2}\) et son inverse \[
F_{p^2}=(I-L_p)^{-1}A_o=\sum_{j\ge0}L_p^jA_o.
\] La somme est finie en chaque point ; sur la norme énergétique, la
série des coûts est majorée par \(\sum_jp^{-j/2}\). Ainsi \[
\frac{\mathcal N_T(A_o)}{1+p^{-1/2}}
\le\mathcal N_T(F_{p^2})
\le\frac{\mathcal N_T(A_o)}{1-p^{-1/2}}.
\] La stabilité est uniforme en \(T\). Elle ne prouve pas que ces normes
restent bornées lorsque \(T\) croît : elle compare les copies à leur
source.

Plus généralement, si \(u=rs^2\) avec \(r\) sans carré, posons
\(A_r=F_r\). Les dilatations commutent et la même preuve donne \[
F_{rs^2}=\prod_{\substack{p\mid s\\p\text{ premier impair}}}
\left(I-\left(\frac rp\right)L_p\right)^{-1}A_r.
\] Cette grammaire d'opérateurs décrit exactement toutes les lignes, y
compris les copies issues de \(r=1\). Les facteurs 2 de \(s\) ne
changent rien, les dénominateurs étant impairs. Cette écriture ne rend
pas les copies indépendantes.

\newpage

\hypertarget{regrouper-la-dette-selon-son-conducteur}{%
\section{3. Regrouper la dette selon son
conducteur}\label{regrouper-la-dette-selon-son-conducteur}}

Pour un entier \(H\ge1\), introduisons l'énergie moyenne \[
Q_H(X)=\sum_{u=1}^H\int_1^X|F_u(x)|^2dx=D_H(X)+R_H(X).
\] La diagonale est positive et satisfait \[
D_H(X)=\sum_{\substack{n\le X\\n\text{ impair}}}\mu(n)^2
\#\{u\le H:(u,n)=1\}\int_n^XB(x/n)^2dx
\le\frac{13}{4}H(X^2-1).
\]

Pour deux indices sans carré, écrire \(n=ga\), \(m=gb\). Alors \(a,b,g\)
sont impairs, sans carré et deux à deux premiers entre eux. Avec
\(q=ab\), \[
\mu(n)\mu(m)=\mu(q),\qquad
\left(\frac un\right)\left(\frac um\right)
=\mathbf1_{(u,g)=1}\left(\frac uq\right).
\] Le facteur commun disparaît du signe de Möbius et devient une
condition de coprimalité. Le signe restant appartient au seul indice
\(q\).

Posons \[
\begin{aligned}
K_X(g;a,b)&=\int_{g\max(a,b)}^X
B\!\left(\frac{x}{ga}\right)B\!\left(\frac{x}{gb}\right)dx,\\
S(H;q,g)&=\sum_{\substack{u\le H\\(u,g)=1}}\left(\frac uq\right).
\end{aligned}
\] Le noyau vaut zéro si sa borne inférieure dépasse \(X\). Le
regroupement exact de toute la dette est \[
\boxed{
R_H(X)=\sum_{\substack{q>1\\q\text{ impair, sans carré}}}\mu(q)
\sum_{a\mid q}
\sum_{\substack{g\le X/\max(a,q/a)\\
g\text{ impair, sans carré}\\(g,q)=1}}
K_X(g;a,q/a)S(H;q,g).}
\] Tous les couples ordonnés sont comptés : la somme sur \(a\mid q\)
inclut les deux orientations. Le cas \(q=1\) redonne exactement \(D_H\),
sans facteur 2 supplémentaire. Les conducteurs \(q>1\) sont au plus
\(X^2\).

Cette écriture conserve les coupures, les poids et les signes. Ce n'est
pas une assertion de positivité terme par terme : \(K_X\), \(S\) et
\(\mu(q)\) peuvent avoir des signes différents.

\newpage

\hypertarget{un-domaine-croissant-contruxf4luxe9-sans-rh}{%
\section{4. Un domaine croissant contrôlé sans
RH}\label{un-domaine-croissant-contruxf4luxe9-sans-rh}}

\textbf{Lemme des petits conducteurs.} Pour \(X\ge1\), \(Q\ge2\) et tout
entier \(H\ge1\), la partie \(1<q\le Q\) de la somme précédente vérifie
\[
\boxed{\begin{aligned}
|R_{1<q\le Q}(X,H)|\le{}&X^2Q(1+\log X)(1+\log Q)^2\\
&\times\left(\frac{13}{2}+\frac54\log Q\right).
\end{aligned}}
\]

\textbf{Preuve, étape 1 : Fourier et coprimalité.} Pour \(q>1\) impair
sans carré, \(\chi_q(v)=(v/q)\) est primitif et non principal. Sa somme
de Gauss a module \(\sqrt q\). La transformation de Fourier finie, la
somme géométrique et \(\sin(\pi j/q)\ge2j/q\) pour \(1\le j\le(q-1)/2\)
donnent \[
\left|\sum_{v\le Y}\chi_q(v)\right|
\le\sqrt q\,h_{(q-1)/2}\le\sqrt q(1+\log q).
\] La borne est uniforme en \(Y\) : les périodes complètes ont une somme
nulle. L'inclusion-exclusion donne exactement \[
S(H;q,g)=\sum_{d\mid g}\mu(d)\left(\frac dq\right)
\sum_{v\le H/d}\left(\frac vq\right),
\quad |S(H;q,g)|\le\tau(g)\sqrt q(1+\log q),
\] où \(\tau(g)\) est le nombre des diviseurs de \(g\).

\textbf{Étape 2 : payer tous les multiplicateurs communs.} Si
\(a\le b\), le changement \(x=gu\) puis la comparaison
\(\sum_{g\le Y}g\tau(g)\le Y^2(1+\log Y)\) entraînent \[
\sum_{g\le X/b}\tau(g)|K_X(g;a,b)|
\le X^2(1+\log X)\int_b^\infty
\frac{|B(u/a)B(u/b)|}{u^2}du.
\] La borne de \(B\), avec \(u=bv\), donne \[
\int_b^\infty\frac{|B(u/a)B(u/b)|}{u^2}du
\le\frac1b\left(\frac{13}{2}+\frac54\log\frac ba\right).
\]

\textbf{Étape 3 : sommer les conducteurs.} Puisque \(ab=q\),
\(\sqrt q/\max(a,b)\le1\) et \(\log(b/a)\le\log q\). Utiliser enfin
\(\sum_{q\le Q}\tau(q)\le Q(1+\log Q)\). Supprimer les restrictions de
coprimalité, de parité et de carré dans ces majorations positives ne
fait qu'augmenter le membre de droite. Cela prouve le lemme.

\textbf{Conséquence.} Choisir \(Q=H\) contrôle la tête à \(HX^2\), avec
des logarithmes. À \(H=X\ge2\), elle est \(O(X^3(1+\log X)^4)\). Le
reste est situé dans \[
\boxed{X<q\le X^2,\qquad g<\sqrt X.}
\] La borne de la tête est absolue. Elle ne constitue pas une loi de
signe globale.

\newpage

\hypertarget{pourquoi-le-grand-crible-usuel-ne-termine-pas}{%
\section{5. Pourquoi le grand crible usuel ne termine
pas}\label{pourquoi-le-grand-crible-usuel-ne-termine-pas}}

Le grand crible quadratique primitif contrôle les lignes impaires sans
carré. Sur un bloc \(n\asymp D\), il donne un coût
\((H+D)\sum_n|a_n|^2\), à un facteur \((HD)^\varepsilon\) près
{[}1,2{]}. Les paramètres carrés nécessaires aux copies n'appartiennent
pas à cette famille.

Pour réintroduire toutes les lignes \(u=rs^2\), on applique ce résultat
aux coefficients \(a_n\mathbf1_{(n,s)=1}\) puis on somme
\(s\le\sqrt H\). Les lignes paires se traitent en tordant les
coefficients par \((2/n)\). On obtient \[
\sum_{u\le H}\left|\sum_{n\asymp D}a_n\left(\frac un\right)\right|^2
\ll_\varepsilon(HD)^\varepsilon
\left(H+D\sqrt H\right)\sum_n|a_n|^2.
\] Le terme supplémentaire vient de la somme des coûts \(D\), pour
chaque \(s\). Avec le poids exact et les bornes du §1, l'intégration
laisse un terme \(X^3\sqrt H\), à des facteurs
\(X^\varepsilon H^\varepsilon\) près. Après division par le nombre des
copies, ce terme rend une borne de taille \(X^3\) : aucun gain fixe
d'exposant n'en résulte.

\textbf{Le terme ne peut pas disparaître pour des coefficients
arbitraires.} Prenons \(H=D\) et des coefficients \(-1\) sur les
premiers \(D/2<n\le D\), zéro ailleurs. Ces coefficients sont
précisément \(\mu(n)\) sur ce support. Si \(D>4\), chaque carré
\(u=s^2\), \(s\le\sqrt D\), est premier à tous ces indices ; sa ligne
vaut donc moins le nombre \(L_D\) de premiers du support. Ces seules
lignes imposent \[
\sum_{u\le D}\left|\sum_n a_n\left(\frac un\right)\right|^2
\ge\lfloor\sqrt D\rfloor L_D^2,
\qquad L_D\asymp\frac D{\log D}.
\] Cela exclut une borne universelle
\(\ll_\varepsilon D^{2+2\varepsilon}\) pour \(\varepsilon<1/4\). Le
masque limité aux premiers diffère du poids exact \(B\) : \textbf{ce
contre-modèle ne réfute pas l'estimation spéciale recherchée}. Il prouve
qu'elle doit utiliser les relations entre premiers et composés, au-delà
du seul fait que les coefficients portent un signe de Möbius.

\textbf{Confrontation des deux analyses.} La construction fournit des
copies exactes et un classement par conducteur. La critique montre que
traiter ces copies par un théorème pour lignes primitives oublie leur
coût. Les deux analyses valident le contrôle des petits conducteurs ;
aucune ne majore encore le reste.

\newpage

\hypertarget{lestimation-restante-et-ce-quelle-donnerait}{%
\section{6. L'estimation restante et ce qu'elle
donnerait}\label{lestimation-restante-et-ce-quelle-donnerait}}

Pour un entier \(X\) grand, poser \(H=X\) et \[
\mathcal P=\{p\text{ premier impair}:\sqrt X/2<p\le\sqrt X\},
\qquad M=|\mathcal P|\asymp\sqrt X/\log X.
\] Le comptage asymptotique des premiers est inconditionnel. Les petites
valeurs où \(\mathcal P\) est vide se traitent séparément sur un
intervalle compact.

Avec \(J_o(X)=\int_1^X|A_o(x)|^2dx\), la restitution exacte et le coût
\(|F_{p^2}(t)|\le5t/2\) donnent \[
\begin{aligned}
M J_o(X)
&\le 2Q_H(X)+\frac{25}{6}X^3\sum_{p\in\mathcal P}\frac1{p^2},\\
J_o(X)&\le\frac{2Q_H(X)}M+\frac{50}{3}X^2.
\end{aligned}
\] La seconde contribution paie la plus petite échelle \(x/p\). Elle est
déjà sous contrôle ; aucune borne de niveau RH n'a été utilisée pour la
payer.

\textbf{Proposition conditionnelle.} Si l'on démontrait, avec \(C,k\)
fixes et indépendants de \(X\), la majoration signée \[
\boxed{R_{q>X}(X;H=X)\le C X^3(1+\log X)^k,}
\] alors le lemme des petits conducteurs et la diagonale donneraient \[
Q_X(X)\ll X^3(1+\log X)^{\max(k,4)},\quad
J_o(X)\ll X^{5/2}(1+\log X)^{\max(k,4)+1}.
\] Puis \(A=A_o-A_o(\cdot/2)\) entraîne \(J_A(X)\le2J_o(X)+4J_o(X/2)\).
La même borne vaut pour \(J\). Le passage des entiers aux réels utilise
la croissance de l'intégrale positive et \(X\le\lceil X\rceil\). La
couverture est donc globale, si l'hypothèse est prouvée.

Pour \(\sigma>3/4\), Cauchy-Schwarz sur \([U,2U]\) donne alors \[
\int_U^{2U}|E(x)|x^{-\sigma-1}dx
\ll U^{3/4-\sigma}(1+\log U)^{(\max(k,4)+1)/2}.
\] La somme dyadique converge. La transformée de Mellin est holomorphe
pour \(\Re s>3/4\), et l'identité \[
-\frac{\zeta'}{\zeta}(s)=\frac{s}{s-1}
+s\int_1^\infty E(x)x^{-s-1}dx
\] y exclut les pôles provenant de zéros. \textbf{La conclusion
conditionnelle est l'absence de zéros pour \(\Re s>3/4\), pas RH.}

Le critère classique \(J(X)=O(X^2)\) vise encore la ligne \(1/2\) et,
par Weil, la positivité sur tous les tests admissibles {[}3,5{]}. Le
demi-plan \(\Re s>3/4\) ne suffirait pas à cette positivité globale.
Aucun gain de cet ordre, ni même le gain conditionnel précédent, n'est
démontré dans cet essai.

\newpage

\hypertarget{calculs-exuxe9cutuxe9s-limites-et-prochaine-attaque}{%
\section{7. Calculs exécutés, limites et prochaine
attaque}\label{calculs-exuxe9cutuxe9s-limites-et-prochaine-attaque}}

Le programme \texttt{calcul\_famille.py} utilise les coefficients
entiers exacts de Möbius et Jacobi. Chaque \(F_u\) est affine sur
\([j,j+1)\) ; son énergie est intégrée par une primitive polynomiale,
sans quadrature. Les coefficients réels sont calculés en double
précision.

\begin{longtable}[]{@{}rrrr@{}}
\toprule\noalign{}
\(X=H\) & \(Q_H/(HX^2)\) & \(D_H/(HX^2)\) & \(R_H/(HX^2)\) \\
\midrule\noalign{}
\endhead
\bottomrule\noalign{}
\endlastfoot
100 & 0,05975949 & 0,05982849 & -0,00006900 \\
300 & 0,05801004 & 0,05861269 & -0,00060265 \\
1000 & 0,05640214 & 0,05766501 & -0,00126287 \\
\end{longtable}

Ces nombres vérifient un calcul fini ; leurs signes n'établissent aucune
loi infinie. Le programme indépendant \texttt{verifier\_identites.py}
recalcule la source en arithmétique décimale à 70 chiffres et vérifie :

\begin{itemize}
\tightlist
\item
  inversion de diviseurs, séparation paire et restitution des copies sur
  sept valeurs réelles et quatre premiers ; résidu maximal
  \(3,04\,10^{-66}\) ;
\item
  regroupement ponctuel par conducteur : résidu \(2\,10^{-68}\) ;
\item
  intégration indépendante à \(X=31\) et \(100\), à partir des valeurs
  et dérivées directes plutôt que des sauts ; différence avec le moteur
  flottant à \(X=100\) inférieure à \(1,1\,10^{-10}\) pour \(Q_H\) et
  \(D_H\) ;
\item
  contre-modèle aux coefficients arbitraires : lignes carrées comptées
  exactement pour cinq tailles, jusqu'à un million.
\end{itemize}

Les diagnostics décimaux ne sont pas des certificats par intervalles. À
\(X=31\), le regroupement intégré donne respectivement environ
\(-36,2662\) et \(-8,5194\) pour les dettes \(q\le X\) et \(q>X\). Cela
ne justifie aucune extrapolation sur la seconde.

\textbf{Prochaine attaque retenue.} Pour \(q>X\), conserver \(\mu(q)\)
dans la somme et appliquer la transformation de Fourier avant de prendre
les valeurs absolues. La transformation quadratique usuelle laisse le
facteur \(\mu(q)\) intact. Il faut une identité supplémentaire qui le
convertisse en une structure estimable, ou une estimation bilinéaire
spéciale de cette dette. La grandeur de sortie exigée est maintenant
écrite au §6 ; aucune de ces deux voies n'est présumée réussir.

\textbf{Garde concernant le préprint sextique {[}4{]}.} Son estimation
moyenne comporte une norme de dérivées du poids lisse. Or
\(B(m)-B(m^-)=\log m+\gamma_E\) : substituer directement \(B(X/n)\)
n'offre pas une constante uniforme. De plus, dans
\(K=\mathbb Q(\sqrt{-3})\), pour \((n,3)=1\), \(N((n))=n^2\) et
\(\mu_K((n))=\mu(n)^2\chi_{-3}(n)\), pas \(\mu(n)\). Un transfert
spectral par \(\zeta_K=\zeta L(\chi_{-3})\) serait possible si le
théorème profond était validé ; sa preuve complète n'est pas auditée
ici. Il n'est utilisé dans aucun lemme de cette note.

\newpage

\hypertarget{ruxe9fuxe9rences-reproduction-et-statut}{%
\section{Références, reproduction et
statut}\label{ruxe9fuxe9rences-reproduction-et-statut}}

\textbf{{[}1{]} D. R. Heath-Brown}, \emph{A mean value estimate for real
character sums}, Acta Arithmetica 72 (1995), 235-275. Notice primaire :
\url{https://www.impan.pl/en/publishing-house/journals-and-series/acta-arithmetica/all/72/3/108575/a-mean-value-estimate-for-real-character-sums}.

\textbf{{[}2{]} L. Goldmakher et B. Louvel}, \emph{A quadratic large
sieve inequality over number fields}, arXiv:1112.1642v2 (2012),
introduction, équation (1.1). \url{https://arxiv.org/pdf/1112.1642}. Le
théorème est appliqué uniquement à sa famille primitive ; le coût de
réintroduction des carrés est démontré séparément au §5.

\textbf{{[}3{]} R. P. Brent, D. J. Platt et T. S. Trudgian}, \emph{The
mean square of the error term in the prime number theorem}, Journal of
Number Theory 238 (2022), 740-762.
\url{https://maths-people.anu.edu.au/~brent/pd/rpb274-final.pdf}. Le
critère énergétique classique demeure une cible, pas une borne
disponible indépendamment de RH. Aucun comportement de plateau de
l'énergie normalisée n'est exigé.

\textbf{{[}4{]} OpenAI}, préprint \emph{The Quasi-Riemann Hypothesis}, 5
octobre 2026, source \texttt{build/paper2.tex}, proposition moyenne et
extraction par puissances.
\url{https://github.com/openai/math/tree/main/preprints/The-Quasi-Riemann-Hypothesis-October-5-2026}.
Source exploratoire, sans recours à ses énoncés profonds. Le dépôt
distingue plusieurs états de vérification :
\url{https://github.com/openai/math}.

\textbf{{[}5{]} M. Suzuki}, \emph{Aspects of the screw function
corresponding to the Riemann zeta function}, arXiv:2206.03682v4 (2023),
théorème 1.3 et §3 sur la positivité de Weil.
\url{https://arxiv.org/pdf/2206.03682}.

\textbf{Provenance locale du raisonnement.} Le dossier V6.54 contient
l'identité des poids premiers, le développement de l'énergie ancrée et
le critère classique. La présente passe démontre les identités de
restitution et le lemme des petits conducteurs dans son propre cadre. La
construction a été confrontée à une analyse critique indépendante ; les
deux ont validé ces preuves auxiliaires. Cette confrontation n'est pas
une validation externe de l'ensemble du programme.

Pour reproduire les diagnostics depuis le dossier :

\begin{Shaded}
\begin{Highlighting}[]
\ExtensionTok{python3}\NormalTok{ verifier\_identites.py }\AttributeTok{{-}{-}output}\NormalTok{ results }\AttributeTok{{-}{-}precision}\NormalTok{ 70}
\ExtensionTok{python3}\NormalTok{ calcul\_famille.py }\AttributeTok{{-}{-}ceilings}\NormalTok{ 100 300 1000}
\end{Highlighting}
\end{Shaded}

Le premier programme utilise seulement la bibliothèque standard Python ;
le second requiert NumPy. Les JSON conservent les résultats et leurs
limites. Le plafond de calcul de la famille est volontairement borné :
les preuves des §§1-6 ne reposent pas sur ce plafond.

\textbf{Synthèse de statut.} Démontré : source exacte, restitution,
stabilité, regroupement, diagonale et tête de conducteurs. Conditionnel
: le transfert d'une borne de la queue vers le demi-plan \(\Re s>3/4\).
Ouvert : cette borne de la queue, la croissance quadratique totale et la
positivité globale de Weil. L'amplification ordinaire ne suffit pas ;
l'amplification exploitant la signature arithmétique de la queue reste à
construire.

\end{document}
